% Generated by scripts/assemble.py; corrections belong in source/corrections.json.
\examhead{117}{福州大学 2025 年数学分析试卷}{113}
% p114-b14; source PDF page 114
\par\Needspace{6\baselineskip}
\blockanchor{p114-b14}
\question{1}(30 分) 求极限.\par
\nopagebreak[4]
% p114-b15; source PDF page 114
\blockanchor{p114-b15}
\displayformula{\lim _ {x \rightarrow 0} \frac {\sqrt {1 + x + x ^ {2}} + \sqrt {1 + x - 2 x ^ {2}} - 2 \sqrt {1 + x}}{e ^ {\cos x - 1} - 1}. \tag {1}}
% p114-b16; source PDF page 114
\blockanchor{p114-b16}
\displayformula{\lim _ {x \to \frac {\pi}{2}} \sum_ {n = 1} ^ {\infty} \frac {\sin (2 n + 1) x}{n ^ {2}}. \tag {2}}
% p114-b17; source PDF page 114
\blockanchor{p114-b17}
\question{2}(20 分) 叙述并证明数列的柯西收敛原理.\par
% p114-b18; source PDF page 114
\blockanchor{p114-b18}
\question{3}(20 分) 设 $f(x)$ 在 $[a, +\infty)$ 上连续, 且 $\lim_{x \to +\infty} [f(x) - 2\sqrt{x}]$ 存在, $f(x)$ 在 $[a, +\infty)$ 上是否一致连续? 如果是或否给出证明, 不一定给出一正一反的例子.\par
% p114-b19; source PDF page 114
\blockanchor{p114-b19}
\question{4}(20 分) 判别数项级数 $\sum_{n=1}^{\infty} \frac{\sin nx}{\sqrt[4]{n}}$ 在 $(0, \pi)$ 上的敛散性, 若收敛指出是绝对收敛还是条件收敛.\par
% p114-b20; source PDF page 114
\blockanchor{p114-b20}
\question{5}(20 分) 证 $\int_{-\infty}^{+\infty} e^{-\frac{x^2}{2}} \, \mathrm{d}x = \sqrt{2\pi}$, 并用该结果计算 $\int_{-\infty}^{+\infty} e^{-\frac{x^2}{2} + x} \, \mathrm{d}x$.\par
% p114-b21; source PDF page 114
\par\Needspace{6\baselineskip}
\blockanchor{p114-b21}
\question{6}(20 分) 计算\par
\nopagebreak[4]
% p114-b22; source PDF page 114
\blockanchor{p114-b22}
\displayformula{\iint_ {\Sigma} \frac {2 x \mathrm{d} y \mathrm{d} z + (2 + z) ^ {2} \mathrm{d} x \mathrm{d} y}{\sqrt {x ^ {2} + y ^ {2} + z ^ {2}}}.}
% p114-b23; source PDF page 114
\blockanchor{p114-b23}
其中 $\Sigma$ 是上半球面 $z = \sqrt{4 - x^2 - y^2}$, 方向取下侧.\par
% p115-b00; source PDF page 115
\par\Needspace{6\baselineskip}
\blockanchor{p115-b00}
\question{7}(20 分) 设 $f(x)$ 在 $[0,1]$ 上连续, 定义 $\alpha$-阶积分算子 $(\alpha > 0)$:\par
\nopagebreak[4]
% p115-b01; source PDF page 115
\blockanchor{p115-b01}
\displayformula{I ^ {\alpha} f (x) = \frac {1}{\Gamma (\alpha)} \int_ {0} ^ {x} (x - s) ^ {\alpha - 1} f (s) \mathrm{d} s.}
% p115-b02; source PDF page 115
\par\Needspace{6\baselineskip}
\blockanchor{p115-b02}
其中 $\Gamma (\alpha)$ 为Gamma函数，利用积分换序原理证：当 $\alpha ,\beta >0$ 时，有\par
\nopagebreak[4]
% p115-b03; source PDF page 115
\blockanchor{p115-b03}
\displayformula{I ^ {\alpha} (I ^ {\beta} f (x)) = I ^ {\alpha + \beta} f (x) = I ^ {\beta} (I ^ {\alpha} f (x)).}
